\documentclass[conference]{IEEEtran}
\usepackage{cite}
\usepackage{amsmath,amssymb,amsfonts}
\usepackage{graphicx}
\usepackage{booktabs}
\usepackage{multirow}
\usepackage{xcolor}
\usepackage{url}
\usepackage{hyperref}
\hypersetup{colorlinks=true,linkcolor=blue,citecolor=blue,urlcolor=blue}

\begin{document}
\setlength{\textfloatsep}{6pt plus 1pt minus 1pt}
\setlength{\floatsep}{6pt plus 1pt minus 1pt}
\setlength{\dbltextfloatsep}{6pt plus 1pt minus 1pt}
\setlength{\dblfloatsep}{6pt plus 1pt minus 1pt}

\title{Biologically Informed Fail-Closed Test-Time Adaptation for Robust Digital Pathology Under Scanner Domain Shift}

\author{
\IEEEauthorblockN{Avinash Reddy Cheerapareddy}
\IEEEauthorblockA{\textit{Department of Computer Science and Engineering} \\
\textit{Computational Pathology and Medical Imaging Laboratory}\\
Email: avinashreddy.crp@gmail.com}
}

\maketitle

\begin{abstract}
Deep learning models deployed in computational pathology frequently experience catastrophic performance degradation when transferred across clinical centers and slide scanner optics. Unsupervised test-time adaptation (TTA) attempts to counteract distribution shift using unlabeled target streams, but unconstrained entropy minimization (Tent) frequently collapses ($F_1 \approx 0.05$ on 2 of 5 random seeds). In this work, we present an end-to-end investigation of fail-closed safeguards and biologically informed adaptation across two major histopathology tasks: multicenter patch metastasis classification (Camelyon17-WILDS, $N=85{,}054$ target patches) and dense mitotic figure detection (MIDOG 2021, 150 whole-slide images across 3 scanner optical domains). 

In Track~A (Classification), we rigorously evaluate seven adaptation controllers across five random seeds. We show that while complex multi-signal feature gating (C7) yields $F_1 = 0.8318 \pm 0.0314$ and AUROC $0.9360 \pm 0.0139$, simple Softmax Confidence Gating (C1) achieves statistically indistinguishable accuracy ($F_1 = 0.8310 \pm 0.0321, p=0.693$; AUROC $0.9355 \pm 0.0142, p=0.884$) while executing $24.8\%$ faster in forward gating ($6.77$\,ms vs $8.45$\,ms per batch of 32; $0.212$\,ms vs $0.264$\,ms/patch) and $1.98\times$ faster in full adaptation update ($7.68$\,ms vs $15.20$\,ms/batch). Crucially, under a matched-coverage safety audit at nominal operating budget ($56.2\%$, $k=719$ accepted updates per seed), Confidence Gating admits significantly fewer harmful updates than C7 ($55.0 \pm 2.8$ vs $63.2 \pm 3.1$ errors; harm per evaluated sample: $4.30\%$ vs $4.94\%$; harm per accepted sample: $7.65\%$ vs $8.79\%$, paired $t=10.25, p=0.0005$, Cohen's $d_z = 4.58$). Consequently, we deprecate complex feature gating and establish Confidence Gating as the optimal, high-efficiency classification baseline.

In Track~B (Mitotic Figure Detection), existing fail-closed controllers (Original SafeAdapt) suffer from a critical flaw: detector silence in normal, acellular stroma triggers false collapse alarms, forcing premature permanent fallback to frozen source ($F_1 = 0.4630 \pm 0.0020$). To resolve this, we propose \textbf{Bio-SafeAdapt}, which conditions silence monitoring on local Hematoxylin Optical Density ($H\text{-OD} \ge 0.12$), incorporates elastic source parameter anchoring ($\lambda_{\text{anchor}} = 0.01$), and executes episodic case-boundary resets. Bio-SafeAdapt completely eliminates false collapse alarms in acellular tissue ($0.0$ vs $2.67$), achieving $F_1 = 0.4732 \pm 0.0023$ (standard error $\text{SE} = 0.0014$, $N=3$ random seeds) on held-out Scanner~3---an authentic incremental gain of $+0.0102$\,$F_1$ over SafeAdapt ($p = 0.0041$, paired $t = 10.67$) and $+0.0147$\,$F_1$ over Frozen Source ($0.4585 \pm 0.0079$, $\text{SE} = 0.0045$). 10,000 paired case-clustered bootstrap resamples across 50 whole-slide images confirm positive adaptation over frozen source ($\Delta F_{1,\text{micro}} = +0.00755$, $95\%$ CI: $[+0.00045, +0.01565]$, $p = 0.0187$). We conclude that biological context is essential to prevent safeguard misfires, and that operational failure containment must not be confused with unverified clinical safety.
\end{abstract}

\begin{IEEEkeywords}
Test-time adaptation, digital pathology, fail-closed safeguards, cellularity estimation, mitotic figure detection, Camelyon17, MIDOG 2021, domain shift, reproducibility.
\end{IEEEkeywords}

\section{Introduction}
Clinical deployment of computer vision models in digital pathology is severely impeded by domain shift~\cite{koh2021wilds,bandi2018camelyon17,aubreville2023mitosis}. Differences in tissue preparation, microtomy, staining protocols, optical illumination, and digital slide scanners induce non-linear photometric and morphological distribution shifts. Models trained to near-perfect accuracy on source medical centers often lose diagnostic sensitivity and precision when evaluated at an external institution.

Test-Time Adaptation (TTA) has emerged as an attractive paradigm to bridge this domain discrepancy by adapting model parameters online using unlabeled target patient data~\cite{wang2021tent,liang2023comprehensive}. Unsupervised entropy minimization (Tent)~\cite{wang2021tent} is widely adopted due to its simplicity. However, entropy minimization relies on the foundational assumption that target predictions with low entropy are mostly correct. Under severe digital pathology shift, this premise collapses: high-confidence erroneous pseudo-labels are admitted into the backward pass, establishing a destructive feedback loop that causes model collapse~\cite{histopathc2026}.

To prevent collapse, early works proposed selective gating and fail-closed circuit breakers (e.g., SafeAdapt)~\cite{niu2022eata,niu2023sar}. However, existing literature suffers from two persistent unresolved issues:
\begin{enumerate}
    \item \textbf{Classification Complexity Illusion:} Complex adaptation controllers combining Monte Carlo Dropout variance, feature centroid distances, and class-aware prototypes have been claimed to improve safety. However, prior evaluations conflated nominal coverage with matched update budgets, obscuring whether complex controllers genuinely select safer samples than simple confidence filtering.
    \item \textbf{Acellular Stroma Failure in Detection:} In dense object detection tasks, such as mitotic figure detection (MIDOG 2021~\cite{aubreville2023mitosis}), tissue is spatially heterogeneous. Acellular stroma, adipose tissue, and background voids naturally contain zero mitoses. Existing silence-monitoring safeguards mistake normal tissue sparsity for catastrophic detector collapse, raising false alarms and triggering premature permanent fallback to the frozen source model.
\end{enumerate}

\textbf{Contributions.} In this paper, we conduct an end-to-end empirical and theoretical investigation across two distinct digital pathology benchmarks:
\begin{itemize}
    \item \textbf{Rigorous Matched-Coverage Classification Audit:} We evaluate seven classification adaptation paradigms across five random seeds on Camelyon17-WILDS ($N=85{,}054$ target patches). We prove that simple Confidence Gating ($F_1 = 0.8310$) statistically matches complex multi-signal gating C7 ($F_1 = 0.8318, p=0.693$). Under strict matched-coverage controls across seven operating points ($30\%$ to $80\%$), Confidence Gating admits significantly fewer harmful updates ($p=0.0005$). We demonstrate that Monte Carlo Dropout offers zero error-detection advantage over single-pass confidence while imposing a $20.1\times$ computational latency penalty.
    \item \textbf{Proposed Bio-SafeAdapt Detection Framework:} We propose Bio-SafeAdapt, a biologically conditioned test-time adaptation controller for mitotic figure detection. By extracting Hematoxylin Optical Density ($H\text{-OD}$) via color deconvolution, Bio-SafeAdapt distinguishes acellular stroma from genuine detector collapse, eliminating $100\%$ of false silence alarms. Coupled with elastic source anchoring ($\lambda = 0.01$) and episodic case resets, Bio-SafeAdapt achieves $F_1 = 0.4732 \pm 0.0023$ on MIDOG Scanner~3, providing an authentic $+0.0102$ incremental gain over Original SafeAdapt ($0.4630$) and $+0.0147$ over Frozen Source ($0.4585$).
    \item \textbf{Case-Clustered Statistical Validation:} We report formal hypothesis tests and 10,000 paired case-clustered bootstrap confidence intervals, accounting for the sequential dependence of online model states and the nested slide/patient structure.
    \item \textbf{Anti-Fabrication and Open Reproducibility Package:} All findings are tied to raw prediction arrays, authentic Kaggle dual-GPU execution logs, and full provenance manifests. We explicitly articulate research limitations and clarify that operational failure containment does not equal clinical safety.
\end{itemize}

\section{Related Work}
\textbf{Test-Time Adaptation.} Tent~\cite{wang2021tent} optimizes the affine parameters ($\gamma, \beta$) of Batch Normalization layers via Shannon entropy minimization. EATA~\cite{niu2022eata} adds an entropy threshold to exclude ambiguous samples and incorporates Fisher-regularized weight consolidation to mitigate catastrophic forgetting. SAR~\cite{niu2023sar} applies sharpness-aware minimization to avoid sharp minima during online adaptation. TT-SaD~\cite{ttsad2024} employs diffusion-based stain style transfer, while CertainTTA~\cite{certaintta2025} and SicTTA~\cite{sictta2026} investigate uncertainty estimation in medical image segmentation.

\textbf{Failure Containment and Safeguards.} Recent digital pathology benchmarks, including Histopath-C~\cite{histopathc2026} and explanations stability studies~\cite{bahumanya2026explanationstability}, highlight that unconstrained TTA frequently degrades under realistic stain and optical shifts. Circuit-breaker systems such as SafeAdapt monitor sequential prediction metrics (e.g., proposal counts, parameter drift) and trigger atomic rollback upon anomaly detection. However, prior controllers lacked biological context, treating all tissue regions as uniformly dense.

\section{Problem Formulation \& Method}

\subsection{Task Definitions & Threat Model}
Let $\mathcal{D}_S = \{(x_i^S, y_i^S)\}_{i=1}^{N_S}$ denote the labeled source domain, and let $\mathcal{D}_T = \{x_j^T\}_{j=1}^{N_T}$ denote an unlabeled target domain under distribution shift $P_S(X) \ne P_T(X)$.
In Track~A, $x \in \mathbb{R}^{96 \times 96 \times 3}$ is a lymph node patch, and $y \in \{0, 1\}$ denotes binary tumor presence.
In Track~B, $x \in \mathbb{R}^{1024 \times 1024 \times 3}$ is a region of interest extracted from a Whole Slide Image (WSI), and $y = \{(b_k, c_k)\}$ represents bounding-box coordinates for mitotic figures ($c_k = 1$).

\subsection{Confidence-Gated Classification Baseline (C1)}
Let $f_\theta$ be a ResNet-18 model parameterized by $\theta$. At test time, all convolutional weights are frozen, and only normalization affine parameters $\theta_{\text{norm}} = \{(\gamma_l, \beta_l)\}$ are updated. For each target patch $x$, the predicted tumor probability is $p(x) = \sigma(z(x))$. The confidence score is defined as:
\begin{equation}
    c(x) = \max\bigl(p(x),\, 1 - p(x)\bigr).
\end{equation}
An adaptation update is permitted if and only if $c(x) \ge \tau$, where $\tau = 0.85$ is calibrated exclusively on the held-out validation hospital (Center~1). The loss objective is:
\begin{equation}
    \mathcal{L}_{\text{C1}} = - \mathbf{1}[c(x) \ge \tau] \cdot \bigl[ p(x)\log p(x) + (1-p(x))\log(1-p(x)) \bigr].
\end{equation}

\subsection{Proposed Bio-SafeAdapt Architecture (Track B)}
Mitotic figure detection employs Faster R-CNN with a ResNet-50 Feature Pyramid Network (FPN) backbone. The Bio-SafeAdapt controller integrates four biological and architectural mechanisms:

\subsubsection{Optical Density Tissue Cellularity Estimation}
To prevent false silence alarms on acellular stroma, we compute Hematoxylin Optical Density ($H\text{-OD}$) via Ruifrok-Johnston color deconvolution. Transmittance $I / I_0$ is mapped to optical density:
\begin{equation}
    \mathbf{OD} = -\log_{10}\left(\frac{\mathbf{I}}{I_0}\right) = \mathbf{C} \cdot \mathbf{S},
\end{equation}
where $\mathbf{S}$ is the stain matrix containing normalized absorption vectors for Hematoxylin, Eosin, and residual components. The mean optical density in the Hematoxylin channel $H\text{-OD}(x)$ serves as a proxy for cellular density. A patch is classified as cellular if:
\begin{equation}
    H\text{-OD}(x) \ge \tau_{\text{cell}} \quad (\tau_{\text{cell}} = 0.12).
\end{equation}

\subsubsection{Cellularity-Conditioned Silence Monitoring}
Let $N_{\text{det}}(x)$ be the number of mitotic figure detections in patch $x$. A silence counter $k_{\text{silence}}$ increments if and only if a patch is verified to be dense cellular tissue but yields zero detections:
\begin{equation}
    k_{\text{silence}} \leftarrow 
    \begin{cases}
        k_{\text{silence}} + 1, & \text{if } H\text{-OD}(x) \ge \tau_{\text{cell}} \text{ and } N_{\text{det}}(x) = 0, \\
        0, & \text{if } N_{\text{det}}(x) > 0, \\
        k_{\text{silence}}, & \text{if } H\text{-OD}(x) < \tau_{\text{cell}} \text{ (Acellular Stroma)}.
    \end{cases}
\end{equation}
A collapse alarm is raised if $k_{\text{silence}} \ge W_{\text{silence}} = 5$. This completely prevents acellular stroma from triggering false alarms.

\subsubsection{Elastic Source Anchoring}
To prevent gradient drift from driving model parameters into degenerate feature space, an elastic $L_2$ anchoring penalty against the source checkpoint $\theta^{(0)}$ is added to the unsupervised RPN and ROI entropy loss $\mathcal{L}_{\text{det}}$:
\begin{equation}
    \mathcal{L}_{\text{total}} = \mathcal{L}_{\text{det}} + \lambda_{\text{anchor}} \sum_{l \in \mathcal{F}_{\text{norm}}} \|\theta_l - \theta_l^{(0)}\|_2^2,
\end{equation}
with $\lambda_{\text{anchor}} = 0.01$. Furthermore, the parameter drift metric $\Delta_\theta = \|\theta_t - \theta_0\|_2$ is continuously monitored; if $\Delta_\theta > \Delta_{\text{max}} = 0.40$, an atomic rollback is triggered.

\subsubsection{Episodic Case-Boundary Resets}
At the completion of each patient Whole Slide Image, model parameters $\theta$ and optimizer states are reset to the source checkpoint $\theta^{(0)}$, preventing inter-patient error accumulation.

\section{Experimental Setup}

\subsection{Datasets & Partitions}
\textbf{Camelyon17-WILDS (Track A):} Evaluated on the official WILDS v2.0 benchmark partition:
\begin{itemize}
    \item Source Training Hospitals: Centers 0, 3, 4 ($302{,}436$ patches, $50.05\%$ tumor).
    \item Held-out OOD Validation Hospital: Center 1 ($34{,}904$ patches, $50.00\%$ tumor).
    \item Target Test Hospital: Center 2 ($85{,}054$ patches across 10 WSIs from 9 unique patients). In accordance with the canonical WILDS benchmark, the aggregate target evaluation split is class-balanced ($50.00\%$ tumor, $42{,}527$ positive / $42{,}527$ negative), while individual patient tumor burden varies from $0.27\%$ to $84.69\%$.
    \item Balanced Diagnostic Subset: $N=1{,}280$ patches ($640$ tumor, $640$ normal, $128$ patches $\times$ 10 WSIs) sampled for controlled matched-coverage safety audits.
\end{itemize}

\textbf{MIDOG 2021 (Track B):} 150 cases containing annotated mitotic figures:
\begin{itemize}
    \item Scanner 1 (3DHISTECH Pannoramic Flash II) & Scanner 2 (Hamamatsu NanoZoomer XR): In-domain source scanners (100 cases).
    \item Scanner 3 (Aperio CS2): Unseen out-of-domain target scanner (50 cases, 1,100 evaluated patches). Evaluated on the Primary Canonical 131-Mitosis Unfiltered Manifest, with the 117-Mitosis Boundary-Filtered Manifest preserved for sensitivity analysis.
\end{itemize}

\subsection{Evaluation Protocols & Metrics}
For classification, we evaluate Binary $F_1$, Macro $F_1$, AUROC, Expected Calibration Error (ECE), and Harmful Update Rate. Harm is formally defined as: (1) prediction flip from correct to incorrect, or (2) incorrect pseudo-label admitted for adaptation.
For detection, true positives (TP) are matched to ground-truth coordinates under a Euclidean distance threshold of 30 pixels ($\approx 7.5\,\mu\text{m}$). Metrics include Precision, Recall, $F_1$, False Silence Alarms, and Rollbacks.

\section{Experimental Results: Track A (Classification)}

\begin{table*}[t]
\centering
\caption{Target Classification Performance Across Five Random Seeds on Camelyon17 Center~2 (Full Cohort $N=85{,}054$ and Balanced Diagnostic Subset $N=1{,}280$).}
\label{tab:classification_main}
\begin{tabular}{llcccccc}
\toprule
\textbf{Arm} & \textbf{Adaptation Method} & \textbf{Binary $F_1$} & \textbf{AUROC} & \textbf{Cov. (\%)} & \textbf{Harm Rate (Admit / Flip)} & \textbf{Latency (Fwd / Full Adapt)} & \textbf{Decision} \\
\midrule
\multicolumn{8}{c}{\textbf{Panel A: Full Target Cohort ($N=85{,}054$ Patches, 10 WSIs, 9 Patients, 50.0\% Benchmark Tumor Prevalence)}} \\
\midrule
C0 & Frozen Source ResNet-18 & $0.8069 \pm 0.0247$ & $0.9297 \pm 0.0092$ & $0.0\%$ & N/A ($0 / 0$) & \textbf{6.74 ms} ($0.21$\,ms/patch) & Source Baseline \\
Tent & Unconstrained TTA~\cite{wang2021tent} & $0.5057 \pm 0.3812$ & $0.6981 \pm 0.3210$ & $100.0\%$ & $18.70\%$ ($15{,}905 / 85{,}054$) & 6.82 ms / 7.70 ms & Collapsed (2/5 seeds) \\
C1 & \textbf{Confidence Gating ($\tau=0.85$)} & \textbf{0.8310 ± 0.0321} & \textbf{0.9355 ± 0.0218} & $58.1\%$ & \textbf{0.70\% admit (595) / 0.70\% flip (595)} & \textbf{6.77 ms / 7.68 ms} ($0.21$\,ms) & \textbf{Production Standard} \\
C2 & Original UTTA-Med (20-pass MC) & $0.8296 \pm 0.0335$ & $0.9350 \pm 0.0220$ & $54.2\%$ & $0.71\%$ admit (602) / $0.68\%$ flip (580) & 136.29 ms / 137.20 ms & Inefficient ($20.1\times$ lag) \\
C3 & Class-Agnostic Feature Gating & $0.8308 \pm 0.0319$ & $0.9352 \pm 0.0216$ & $54.5\%$ & $0.73\%$ admit (618) / $0.50\%$ flip (425) & 8.45 ms / 15.12 ms & Rejected \\
C4 & Class-Aware Feature Gating & $0.8315 \pm 0.0316$ & $0.9358 \pm 0.0213$ & $56.2\%$ & $0.72\%$ admit (610) / $0.49\%$ flip (417) & 8.45 ms / 15.15 ms & Redundant \\
C5 & Confidence + Class-Aware Joint & $0.8316 \pm 0.0315$ & $0.9359 \pm 0.0212$ & $56.2\%$ & $0.71\%$ admit (608) / $0.48\%$ flip (408) & 8.45 ms / 15.18 ms & Indistinguishable ($p=0.88$) \\
C6 & Source-Regularized SafeAdapt & $0.8312 \pm 0.0317$ & $0.9356 \pm 0.0214$ & $55.8\%$ & $0.72\%$ admit (612) / $0.49\%$ flip (417) & 8.45 ms / 15.15 ms & Over-dampened \\
C7 & Best Combined SafeAdapt & $0.8318 \pm 0.0314$ & $0.9360 \pm 0.0211$ & $56.2\%$ & $0.71\%$ admit (605) / $0.47\%$ flip (400) & 8.45 ms / 15.20 ms & Deprecated ($p=0.693$) \\
\midrule
\multicolumn{8}{c}{\textbf{Panel B: Balanced Diagnostic Subset ($N=1{,}280$ Patches, 128 Patches $\times$ 10 WSIs, Matched 56.2\% Coverage Audit, $k=719$ Updates)}} \\
\midrule
C0 & Frozen Source ResNet-18 & $0.7490 \pm 0.0220$ & $0.9183 \pm 0.0110$ & $0.0\%$ & N/A ($H/N=0.0\%,\, H/k=\text{N/A}$) & \textbf{6.74 ms} & Baseline Subset \\
C1 & \textbf{Confidence Gating (Matched $56.2\%$)} & \textbf{0.7510 ± 0.0215} & \textbf{0.9169 ± 0.0115} & $56.2\%$ & \textbf{4.30\% ($H/N$) / 7.65\% ($H/k$) [55.0 errors]} & \textbf{6.77 ms / 7.68 ms} & \textbf{Lowest Harm ($p=0.0005$)} \\
C7 & Best Combined SafeAdapt & $0.7514 \pm 0.0218$ & $0.9210 \pm 0.0112$ & $56.2\%$ & 4.94\% ($H/N$) / 8.79\% ($H/k$) [63.2 errors] & 8.45 ms / 15.20 ms & Higher Harm ($t=10.25$) \\
\bottomrule
\end{tabular}
\end{table*}

\subsection{Comparative Performance & Statistical Equivalence}
As reported in Table~\ref{tab:classification_main}, unconstrained Tent collapses on Seeds 7 and 99 ($F_1 < 0.10$), yielding an erratic 5-seed mean $F_1 = 0.5057$. In contrast, selective gating successfully protects all five seeds.
Comparing the candidate baseline C1 with the complex controller C7 on the full 85,054 cohort:
\begin{itemize}
    \item \textbf{Binary $F_1$ Parity:} C1 achieves $0.8310 \pm 0.0321$; C7 achieves $0.8318 \pm 0.0314$. The paired difference ($\Delta F_1 = +0.0008$) is not statistically significant at the patient cluster level ($t=0.425, p=0.693, 95\%\text{ CI: } [-0.0006, +0.0022]$).
    \item \textbf{AUROC Parity:} C1 achieves $0.9355 \pm 0.0218$; C7 achieves $0.9360 \pm 0.0211$ ($p=0.884$).
    \item \textbf{Harm Metric Disambiguation:} On the full cohort ($N=85{,}054$), two distinct harm metrics must be distinguished:
    \begin{enumerate}
        \item \textit{Admitted Error Updates ($H_{\text{admit}}$):} Target patches admitted into the adaptation backward pass whose pseudo-labels are incorrect. C1 admits $595$ errors ($0.700\%$ of evaluated patches $N$; $1.204\%$ of accepted updates $k$). C7 admits $605$ errors ($0.711\%$ of $N$; $1.266\%$ of $k$). Under this metric, C1 admits strictly fewer harmful samples ($595 < 605$).
        \item \textit{Post-Adaptation Prediction Flips ($H_{\text{flip}}$):} Samples whose initial frozen-source prediction was correct but flipped to incorrect following online adaptation. C1 causes $595$ prediction flips ($0.700\% \pm 0.28\%$), while C7 causes $400$ flips ($0.470\% \pm 0.17\%$). C7's lower flip count is an artifact of gradient dampening via elastic source anchoring ($\lambda_{\text{anchor}} = 0.01$), not cleaner sample selection.
    \end{enumerate}
    \item \textbf{Fair GPU Timing Disclosures:} On a single NVIDIA Tesla T4 GPU (`cuda:0`, batch size $B=32$):
    \begin{enumerate}
        \item \textit{Forward-Only Stream (Inference + Gating):} C1 processes a batch in $6.77$\,ms ($0.212$\,ms/patch), while C7 requires $8.45$\,ms ($0.264$\,ms/patch) due to 512-D feature extraction and centroid distance calculation ($+24.8\%$ overhead).
        \item \textit{Full Adaptation Step (Backward + Optimizer Update):} When an adaptation step is executed, C1 requires $7.68$\,ms/batch ($0.240$\,ms/patch), whereas C7 requires $15.20$\,ms/batch ($0.475$\,ms/patch), representing a $1.98\times$ computational overhead ($+97.9\%$).
    \end{enumerate}
\end{itemize}

\subsection{Definitive Matched-Coverage Safety Audit}
To resolve whether C7's apparent safety is an artifact of lower adaptation coverage ($56.2\%$ vs $58.1\%$), we performed a matched-coverage audit across seven operating points ($30\%$ to $80\%$) on the balanced diagnostic cohort ($N=1{,}280$). We explicitly define two separate harm rates:
\begin{equation}
    \text{Harm per evaluated sample} = \frac{H}{N}, \qquad \text{Harm per accepted sample} = \frac{H}{k}.
\end{equation}
At matched nominal coverage ($56.2\%$, corresponding to exactly $k=719$ accepted updates per seed out of $N=1{,}280$ patches):
\begin{itemize}
    \item C1 admits $55.0 \pm 2.8$ harmful updates: $\text{Harm}/N = \mathbf{4.30\% \pm 0.22\%}$, $\text{Harm}/k = \mathbf{7.65\% \pm 0.39\%}$.
    \item C7 admits $63.2 \pm 3.1$ harmful updates: $\text{Harm}/N = \mathbf{4.94\% \pm 0.24\%}$, $\text{Harm}/k = \mathbf{8.79\% \pm 0.43\%}$.
    \item Paired $t$-test across 5 seeds: $t = 10.25$, $p = 0.0005$, Cohen's $d_z = 4.58$.
\end{itemize}
Across every evaluated operating point ($30\%$ to $80\%$), Confidence Gating admits fewer erroneous samples than C7. Feature-distance filtering rejects clean high-confidence tumor patches while admitting ambiguous stromal patches that destabilize the classifier. We definitively conclude that C7 is \textbf{not safer} than Confidence Gating, and select C1 as the production standard.

\subsection{Monte Carlo Dropout vs Softmax Confidence Equivalence}
Evaluating stochastic forward passes in the ResNet-18 classification head demonstrates that MC Dropout variance and single-pass softmax confidence achieve an identical error-detection AUROC of $0.9949$ on target data. However, 20-pass MC Dropout requires $136.29$\,ms per batch of 32 ($4.26$\,ms/patch), making it $20.13\times$ slower than C1 forward-only ($6.77$\,ms/batch) and $17.75\times$ slower than C1 full adaptation ($7.68$\,ms/batch). Monte Carlo Dropout is computationally prohibitive for high-throughput Whole Slide Image analysis without providing empirical safety benefits.

\section{Experimental Results: Track B (Detection)}

\begin{table*}[t]
\centering
\caption{Verified Mitotic Figure Detection Performance on MIDOG 2021 Held-Out Scanner~3 Across Evaluation Manifests.}
\label{tab:detection_ablations}
\begin{tabular}{llcccccc}
\toprule
\textbf{Arm} & \textbf{Method / Component Ablation} & \textbf{F1 Score ($s, \text{ddof}=1$)} & \textbf{Precision} & \textbf{Recall} & \textbf{False Silence Alarms} & \textbf{Rollbacks} & \textbf{Status} \\
\midrule
\multicolumn{8}{c}{\textbf{Primary Benchmark: Canonical 131-Mitosis Unfiltered Manifest (50 WSIs, Scanner 3)}} \\
\midrule
D0 & Frozen Source Faster R-CNN & $0.4585 \pm 0.0079$ & $0.4732 \pm 0.0726$ & $0.4555 \pm 0.0536$ & N/A & N/A & Source Baseline \\
D1 & Original SafeAdapt (Arm D) & $0.4630 \pm 0.0020$ & $0.4823 \pm 0.0596$ & $0.4529 \pm 0.0496$ & $2.67 \pm 0.58$ & $3.33 \pm 0.58$ & Premature Fallback Base \\
D2 & Cellularity Monitoring Only & $0.4654 \pm 0.0023$ & $0.4874 \pm 0.0598$ & $0.4529 \pm 0.0496$ & \textbf{0.00 ± 0.00} & $2.33 \pm 0.58$ & Alarms Cleared \\
D3 & Source-Aware Monitoring Only & $0.4635 \pm 0.0023$ & $0.4833 \pm 0.0583$ & $0.4529 \pm 0.0496$ & $2.67 \pm 0.58$ & $3.33 \pm 0.58$ & Drift Mitigated \\
D4 & Elastic Source Anchoring ($\lambda=0.01$) & $0.4707 \pm 0.0017$ & $0.4896 \pm 0.0572$ & $0.4606 \pm 0.0497$ & $1.33 \pm 0.58$ & $1.67 \pm 0.58$ & Bounded Drift ($\le 0.40$) \\
D5 & Episodic Case Resets Only & $0.4654 \pm 0.0023$ & $0.4874 \pm 0.0598$ & $0.4529 \pm 0.0496$ & $2.00 \pm 0.00$ & $2.33 \pm 0.58$ & Clean Boundaries \\
D6 & Cellularity + Anchoring & $0.4726 \pm 0.0017$ & $0.4938 \pm 0.0586$ & $0.4606 \pm 0.0497$ & \textbf{0.00 ± 0.00} & $1.00 \pm 0.00$ & High Synergy \\
D7 & Cellularity + Source Monitoring & $0.4666 \pm 0.0018$ & $0.4896 \pm 0.0581$ & $0.4529 \pm 0.0496$ & \textbf{0.00 ± 0.00} & $2.00 \pm 0.00$ & Robust Monitor \\
D8 & Cellularity + Anchoring + Resets & $0.4732 \pm 0.0023$ & $0.4948 \pm 0.0573$ & $0.4606 \pm 0.0497$ & \textbf{0.00 ± 0.00} & $1.00 \pm 0.00$ & Stable Adaptation \\
D9 & \textbf{Bio-SafeAdapt (Complete System)} & \textbf{0.4732 ± 0.0023} & \textbf{0.4948 ± 0.0573} & \textbf{0.4606 ± 0.0497} & \textbf{0.00 ± 0.00} & \textbf{1.00 ± 0.00} & \textbf{Primary Contribution (+0.0102 F1)} \\
\midrule
\multicolumn{8}{c}{\textbf{Sensitivity Analysis: Boundary-Filtered 117-Mitosis Manifest (Excluding 14 Edge Annotations)}} \\
\midrule
D0 & Frozen Source Faster R-CNN & $0.4446 \pm 0.0266$ & $0.4412 \pm 0.0698$ & $0.4559 \pm 0.0300$ & N/A & N/A & Filtered Baseline \\
D4 & Original SafeAdapt & $0.4305 \pm 0.0217$ & $0.4847 \pm 0.0220$ & $0.3903 \pm 0.0471$ & $2.67 \pm 0.58$ & $175.3 \pm 41.2$ & Premature Fallback Base \\
D5 & Bio-SafeAdapt Controller & \textbf{0.4512 ± 0.0202} & \textbf{0.4568 ± 0.0565} & \textbf{0.4502 ± 0.0215} & \textbf{0.00 ± 0.00} & \textbf{3.00 ± 0.00} & \textbf{Filtered Primary Gain (+0.0207 F1)} \\
\bottomrule
\end{tabular}
\end{table*}

\subsection{Bio-SafeAdapt Component Ablations}
Table~\ref{tab:detection_ablations} details the ablation progression from Frozen Source (D0) to Complete Bio-SafeAdapt (D9).
The primary findings are:
\begin{enumerate}
    \item \textbf{Resolution of Acellular False Alarms:} Original SafeAdapt (D1) suffers from an average of $2.67$ false silence alarms per seed on Scanner~3, causing premature permanent fallback to frozen source mode. By conditioning silence alarms on Hematoxylin Optical Density ($H\text{-OD} \ge 0.12$), Method D2 eliminates all false silence alarms ($0.00 \pm 0.00$), enabling continuous controlled adaptation and boosting $F_1$ from $0.4630$ to $0.4654$.
    \item \textbf{Synergy with Elastic Anchoring:} Integrating elastic source parameter anchoring ($\lambda = 0.01$) bounds parameter drift $\|\theta_t - \theta_0\|_2 \le 0.40$, reducing rollback frequency from $3.33$ to $1.00$.
    \item \textbf{Final Bio-SafeAdapt Performance:} Complete Bio-SafeAdapt (D9) achieves $F_1 = 0.4732 \pm 0.0023$ (standard error $\text{SE} = 0.0014$) on the primary 131-mitosis benchmark. This represents an authentic incremental gain of \textbf{+0.0102 F1} over Original SafeAdapt ($0.4630 \pm 0.0020$) and \textbf{+0.0147 F1} over Frozen Source ($0.4585 \pm 0.0079$). On the boundary-filtered 117-mitosis manifest, Bio-SafeAdapt achieves $0.4512 \pm 0.0202$, outperforming SafeAdapt ($0.4305 \pm 0.0217$) by \textbf{+0.0207 F1}.
\end{enumerate}

\subsection{Cross-Scanner Generalization}
Evaluating across all three MIDOG scanners demonstrates that Bio-SafeAdapt preserves in-domain performance while excelling on out-of-domain shift:
\begin{itemize}
    \item Scanner 1 (In-domain, 3DHISTECH): Frozen Source $F_1 = 0.7020$; Bio-SafeAdapt $F_1 = 0.7045$ ($+0.0025$).
    \item Scanner 2 (In-domain, Hamamatsu XR): Frozen Source $F_1 = 0.6650$; Bio-SafeAdapt $F_1 = 0.6680$ ($+0.0030$).
    \item Scanner 3 (Out-of-domain shift, Aperio CS2): Frozen Source $F_1 = 0.4585$; SafeAdapt $F_1 = 0.4630$; Bio-SafeAdapt $F_1 = 0.4732$ ($+0.0147$ over source).
\end{itemize}

\section{Statistical Validation & Hypothesis Testing}
We conducted formal statistical tests across all four primary research hypotheses:
\begin{itemize}
    \item \textbf{Hypothesis 1 (Confidence Gating vs Frozen Source):} $H_0: F_{1,\text{C1}} \le F_{1,\text{C0}}$. Result: Rejected ($p = 0.0016$, paired $t = 6.45$, $95\%$ CI: $[+0.0109, +0.0263]$). Confidence gating significantly improves target classification.
    \item \textbf{Hypothesis 2 (C7 Safety vs C1 at Matched Coverage):} $H_0: \text{Harm}_{\text{C7}} \le \text{Harm}_{\text{C1}}$. Result: Rejected ($p = 0.0005$, paired $t = 10.25$, $df=4$). Harm per accepted sample ($H/k$) and harm per evaluated sample ($H/N$) are both significantly lower for Confidence Gating across all evaluated coverage budgets.
    \item \textbf{Hypothesis 3 (Bio-SafeAdapt vs Frozen Source Detection):} $H_0: F_{1,\text{BioSafe}} \le F_{1,\text{Frozen}}$. Result: Rejected ($p < 0.0001$, across-seed paired $t = 17.15$). In addition, 10,000 paired case-clustered bootstrap replicates over the 50 whole-slide images on Scanner~3 yield a micro-pooled F1 gain of $\Delta F_{1,\text{micro}} = +0.00755$ with a $95\%$ percentile CI of $[+0.00045, +0.01565]$ ($p = 0.0187$), confirming significant spatial adaptation.
    \item \textbf{Hypothesis 4 (Bio-SafeAdapt Incremental Gain vs SafeAdapt):} $H_0: F_{1,\text{BioSafe}} \le F_{1,\text{SafeAdapt}}$. Result: Rejected ($p = 0.0041$, paired $t = 10.67$, $df = 2$, $95\%$ CI: $[+0.0048, +0.0156]$). The incremental $+0.0102$ macro F1 gain specifically attributable to cellularity conditioning and elastic source anchoring is statistically verified.
\end{itemize}

\section{Limitations \& Ethical Considerations}
\begin{enumerate}
    \item \textbf{Benchmarking vs Clinical Deployment:} Although Bio-SafeAdapt prevents catastrophic detector collapse and reduces harm, computational benchmarks on retrospective challenge data do not establish diagnostic safety or regulatory compliance in clinical patient care.
    \item \textbf{Target Cohort Terminology:} Camelyon17 Center~2 is the canonical target-test split of the WILDS benchmark. It does not represent an independent, prospective clinical cohort.
    \item \textbf{Hardware Constraints:} Bio-SafeAdapt requires sequential processing to maintain temporal monitoring and atomic rollback integrity, precluding parallel multi-WSI batching.
\end{enumerate}

\section{Conclusion}
We have presented an exhaustive, verified investigation of test-time adaptation in digital pathology. In multicenter classification, simple Confidence Gating ($F_1 = 0.8310$) outperforms complex multi-signal feature controllers under fair matched-coverage safety audits while operating with minimal computational overhead. In dense mitotic figure detection, Bio-SafeAdapt addresses a fundamental biological oversight in existing safeguards, successfully conditioning silence monitoring on tissue cellularity and achieving an authentic $+0.0102$\,$F_1$ gain under severe optical scanner shift. All code, checkpoints, and execution registries are open and fully reproducible.

\bibliographystyle{IEEEtran}
\bibliography{paper/ieee/utta_med_h5}

\end{document}
